\documentclass[aspectratio=169]{beamer}
\usetheme{default}
\usecolortheme{dove}
\usefonttheme{serif}
\setbeamertemplate{navigation symbols}{}
\setbeamertemplate{footline}[frame number]
\beamertemplatenavigationsymbolsempty

\usepackage{amsmath,amssymb}
\usepackage{graphicx}

\graphicspath{{figs/}}

\definecolor{key}{rgb}{0.10,0.35,0.65}
\setbeamercolor{frametitle}{fg=key}

\title{Stellarator Neutronics Check In Meeting}
\author{Thaddaeus Kiker}
\date{2026-08-19}

\begin{document}

\begin{frame}[plain]\titlepage\end{frame}

\begin{frame}{The Device}
  \begin{columns}[c]
    \begin{column}{0.7\textwidth}
      \includegraphics[width=\textwidth]{device_wireframe}
    \end{column}
    \begin{column}{0.3\textwidth}
      \begin{itemize}
        \item Quasi-helical (Landreman--Paul) stellarator, nfp = 4.
      \end{itemize}
    \end{column}
  \end{columns}
\end{frame}

\begin{frame}{Surface Source Decoupling}
  \begin{itemize}
    \item Goal: run the expensive plasma transport once, then reuse it across many shield designs.
    \item I freeze the neutron field at the first wall using surface-source read/write (SSR), then replay those neutrons through each blanket and shield variant out to the magnets.
    \item Freezing in front of the blanket keeps both breeder and shield thickness free to vary.
  \end{itemize}
\end{frame}

\begin{frame}{The Double-Counting Mistake, and the Fix}
  \begin{itemize}
    \item The freeze records a neutron every time it crosses the first wall.
    \item The wall wraps a torus, so a neutron can scatter, cross back through the hole, and be recorded again --- about 2.4 records per neutron.
    \item 58\% of the frozen neutrons were these re-crossings, not fresh plasma neutrons.
    \item Replaying them as independent particles double-counted, and skewed the coil-load pattern.
    \item I found it by energy: fresh plasma neutrons are still at 14.06 MeV, re-crossings are degraded below it.
    \item Fix: keep only the fresh uncollided neutrons and leave the void empty, so re-crossing happens naturally during the replay.
    \item Partial fix: the peak-load coils now match the full run; the far coils still come out low.
  \end{itemize}
\end{frame}

\begin{frame}{Full Run vs SSR}
  \begin{columns}[c]
    \begin{column}{0.62\textwidth}
      \includegraphics[width=\textwidth]{ci_fluxbars}
    \end{column}
    \begin{column}{0.38\textwidth}
      \begin{itemize}
        \item Each series is scaled to its own peak coil, so we compare the load shape.
        \item Inboard peak-load coils, 16 to 21, sit behind the highest-current wall; the two agree there.
        \item Outboard far coils, 26 to 30, are fed by indirect scattered paths; SSR comes out low.
      \end{itemize}
    \end{column}
  \end{columns}
\end{frame}

\begin{frame}{The Speed-Up Question}
  \begin{itemize}
    \item One SSR replay ran about 7h48m against the full run's 8h48m --- roughly a 10 to 13\% saving at equal particle count (24M).
    \item I do have the infrastructure to mesh the neutron field at any resolution we want --- for example on the neutron-exit layer coming out of the shield onto the magnets.
  \end{itemize}
\end{frame}

\begin{frame}{Next Steps}
  \begin{itemize}
    \item Localize the far-coil residual: tally the void-crossing current in the full run versus the replay.
    \item Test freezing at the neutron-exit layer, where the magnet load is captured directly and should be exact for every coil.
    \item Quantify the correlated-sampling saving: replay through modified builds against their finite-difference truth.
    \item Settle the freeze-surface choice: the first wall keeps the breeder trade, deeper is faster but fixes the breeder.
  \end{itemize}
\end{frame}

\end{document}
